\documentclass[11pt,a4paper]{article}
\usepackage[margin=2.5cm]{geometry}
\usepackage{times}
\usepackage{latexsym}
\usepackage{graphicx}
\usepackage{amsmath}
\usepackage{amssymb}
\usepackage{booktabs}
\usepackage{multirow}
\usepackage{xcolor}
\usepackage{hyperref}
\usepackage{url}
\usepackage{microtype}
\usepackage{array}
\usepackage{tabularx}
\usepackage{natbib}
\usepackage{enumitem}

\newcommand{\method}{CaCLR}
\newcommand{\rmark}{$\rho$}

\title{Cross-Lingual Correctness Probe Transfer:\\Negative Results, Methodological Warnings, and Conditional Solutions}

\author{Anonymous}

\begin{document}
\maketitle

\begin{abstract}
Probe-based calibration methods extract internal correctness directions from LLM representations, but all published work is English-only.
We present the first large-scale study of cross-lingual correctness probe transfer, testing three architectures (Qwen3-235B MoE, Llama~4 Scout MoE, Gemma~4 31B Dense) across 20 typologically diverse languages.
Our findings are predominantly negative but informative.
(1)~\textbf{Control task warning}: extending Hewitt \& Liang selectivity checks to cross-lingual probing reveals that 35\% of language--layer combinations are capacity artifacts (LOW-resource: 42\%, MID: 31\%), with selectivity strongly correlated with zero-shot AUROC ($\rho{=}0.65$). This implies many published positive cross-lingual probing results may reflect artifacts rather than genuine signal transfer.
(2)~\textbf{Mean-shift impossibility}: under StandardScaler + logistic regression, mean-shift produces $\Delta$AUROC~$\equiv$~0 by rank-invariance, confirmed across all 60 model--language pairs.
(3)~\textbf{Distributional matching path closed}: CaCLR (MMD), CORAL, and high-dimensional MMD (PCA $d{=}50,100$) all degrade performance; every added loss component is net-negative in ablation.
(4)~\textbf{Linear projection guidance}: Procrustes alignment with ${\geq}100$ parallel pairs is the only consistently beneficial label-free method, but its effect reverses for languages with zero-shot AUROC~$>$~0.62, precluding a single actionable threshold.
We release all code and negative results to prevent redundant exploration of closed paths.
\end{abstract}

%-------------------------------------------------------------
\section{Introduction}
\label{sec:intro}

Recent work has demonstrated that large language models encode linear correctness directions in their internal representations~\cite{miao2026coral,burns2023discovering}.
Probes trained on English representations achieve AUROC $>$~0.75 for predicting answer correctness, enabling representation-level calibration that is complementary to output-based methods.
However, all probe-based calibration results to date are English-only, leaving a critical question unanswered: \emph{do these correctness directions transfer across languages?}

This question matters because non-English languages exhibit substantially worse calibration~\cite{ahuja2024globalmmlu}, and the most promising solutions---per-language supervised probes or ensemble methods like LACE---require target-language labels that are expensive for low-resource settings.
If English correctness directions transfer, they offer a zero-label calibration path.

A pilot study on Qwen2.5-7B and Llama-3.1-8B (5 languages each) suggested promising patterns: a local Pareto tradeoff at deepest layers, a Simpson's Paradox in CKA--transfer correlations, and a natural distributional matching approach (CaCLR) to test.
We scaled this to three production architectures, 20 languages, and systematic ablations.

\textbf{The results are predominantly negative.}
The pilot patterns do not generalize.
MoE architectures show no Pareto frontier; the only dense model (Gemma~4 31B) shows marginal effects far weaker than the pilot.
The Simpson's Paradox in CKA is a small-sample artifact.
CaCLR and all distributional matching variants fail.
Difficulty--capability decomposition is indistinguishable from random projection.

However, we contribute four findings of lasting value:
\begin{enumerate}
\item \textbf{Control task methodology warning} (\S\ref{sec:control}): 35\% of cross-lingual probe results are capacity artifacts, with implications for the entire probing literature.
\item \textbf{Mean-shift impossibility} (\S\ref{sec:meanshift}): a formal proof that mean-shift cannot improve AUROC under standard probe pipelines.
\item \textbf{Distributional matching impossibility} (\S\ref{sec:caclr}): systematic closure of MMD, CORAL, and high-dimensional variants.
\item \textbf{Conditional linear projection guidance} (\S\ref{sec:linproj}): the only viable label-free path, with resource-tier-dependent sample efficiency.
\end{enumerate}

%-------------------------------------------------------------
\section{Related Work}
\label{sec:related}

\paragraph{Probe-based calibration.}
Burns et al. extract contrast-consistent directions; Miao et al. train correctness probes with cross-task transfer (CORAL); Miao \& Ungar demonstrate correctness--confidence orthogonality and probe-based steering.
All work exclusively in English.

\paragraph{Cross-lingual probing.}
Xiao et al. transfer knowledge-boundary probes (input-level answerability) across 8 languages.
Critically, answerability is output-independent while correctness depends on the model's actual answer---a fundamentally different transfer problem.
Civelli et al. transfer difficulty probes cross-lingually; difficulty serves as our ablation control.

\paragraph{Domain adaptation theory.}
Ben-David et al. bound target-domain error via source error plus $\mathcal{H}$-divergence.
Zhao et al. prove that marginal distribution matching fails when conditional distributions $P(Y|X)$ differ across domains---directly predicting our CaCLR negative result.

\paragraph{Cross-lingual alignment.}
MUSE and VecMap learn linear projections between embedding spaces.
Our Procrustes alignment is conceptually analogous at the sentence-level correctness representation.

\paragraph{Probing methodology.}
Hewitt \& Liang propose control tasks measuring probe selectivity; Voita \& Titov use MDL probes.
We extend selectivity checks to the cross-lingual setting for the first time.

%-------------------------------------------------------------
\section{Experimental Setup}
\label{sec:setup}

\subsection{Models}
We probe three architectures:
\textbf{Qwen3-235B-A22B} (94 layers, MoE, INT4 quantized, 24 probed layers every 4th);
\textbf{Llama~4 Scout} (48 layers, 17B-16E MoE, W4A16, 12 probed layers every 4th);
\textbf{Gemma~4 31B} (60 layers, dense, BF16 unquantized, 16 probed layers every 4th + final).
Gemma serves as a quantization-free dense control.

\subsection{Languages and Data}
We evaluate 20 typologically diverse languages from GlobalMMLU, grouped into three resource tiers:
\textbf{HIGH} (7): zh, es, fr, de, ja, pt, it;
\textbf{MID} (7): ar, hi, ru, ko, tr, vi, th;
\textbf{LOW} (6): sw, bn, am, yo, te, my.
Probes are trained on the English portion (70\% stratified split for training, 30\% for evaluation).
Additional benchmarks (Belebele, MMLU-ProX) are evaluated in Appendix~\ref{app:benchmarks}.

\subsection{Transfer Conditions}
\label{sec:conditions}
Five conditions are compared:
(1)~\textbf{Zero-shot}: apply English-trained probe directly;
(2)~\textbf{Mean-shift}: subtract mean activation difference between English and target;
(3)~\textbf{Linear projection} (Procrustes, 200 parallel pairs);
(4)~\textbf{Few-shot} (25 or 100 target-language labels);
(5)~\textbf{CaCLR}: joint calibration + distributional matching (details in \S\ref{sec:caclr}).

\subsection{Probes and Metrics}
Logistic regression probes (justified in Appendix~\ref{app:probe}); AUROC for discrimination, ECE for calibration.
A language ``passes'' zero-shot transfer if AUROC~$>$~0.58.

%-------------------------------------------------------------
\section{Mean-Shift Impossibility}
\label{sec:meanshift}

Under StandardScaler preprocessing followed by logistic regression, mean-shift adds a constant offset $c$ to all probe logits:
$\hat{y}_{\text{shifted}}(x) = w^\top(x + \delta) + b = \hat{y}(x) + w^\top\delta$.
Since AUROC depends only on ranking, and a constant offset preserves ranking, $\Delta$AUROC~$\equiv$~0 \emph{mathematically}.

All 60 model--language combinations confirm $\Delta$AUROC~$=$~0.000 (Tables~\ref{tab:qwen_transfer},~\ref{tab:llama_transfer},~\ref{tab:gemma_summary}).
This distinguishes our pipeline from Xiao et al., where mean-shift operates under different preprocessing.

%-------------------------------------------------------------
\section{Pareto Frontier Analysis}
\label{sec:pareto}

The pilot study on 7--8B dense models suggested a depth--transfer Pareto tradeoff.
We test this across three full-scale architectures.

\subsection{MoE Architectures: No Pareto}

\begin{table}[t]
\centering
\small
\caption{Qwen3-235B: per-language depth--transfer Spearman $\rho$. Only 1/20 languages reaches significance.}
\label{tab:qwen_pareto}
\begin{tabular}{llcccc}
\toprule
Lang & Tier & ZS AUROC & $\rho$ & $p$ & Pattern \\
\midrule
zh & HIGH & .668 & $-.31$ & .14 & weak neg \\
es & HIGH & .675 & $-.25$ & .24 & n.s. \\
fr & HIGH & .658 & $-.28$ & .18 & n.s. \\
de & HIGH & .649 & $-.19$ & .38 & n.s. \\
ja & HIGH & .588 & $+.18$ & .40 & U-shape \\
pt & HIGH & .662 & $-.22$ & .30 & n.s. \\
it & HIGH & .645 & $-.16$ & .46 & n.s. \\
\midrule
ar & MID & .608 & $+.05$ & .82 & flat \\
hi & MID & .567 & $-.41$ & \textbf{.046} & sig.\ neg \\
ru & MID & .621 & $-.14$ & .52 & n.s. \\
ko & MID & .591 & $+.25$ & .24 & inverted \\
tr & MID & .548 & $-.08$ & .71 & n.s. \\
vi & MID & .562 & $+.09$ & .68 & flat \\
th & MID & .528 & $-.11$ & .61 & n.s. \\
\midrule
sw & LOW & .518 & $-.35$ & .09 & marginal \\
bn & LOW & .524 & $-.12$ & .58 & n.s. \\
am & LOW & .486 & $+.14$ & .52 & non-mono \\
yo & LOW & .475 & $-.02$ & .93 & flat \\
te & LOW & .502 & $-.18$ & .40 & n.s. \\
my & LOW & .468 & $+.10$ & .64 & non-mono \\
\bottomrule
\end{tabular}
\end{table}

\textbf{Qwen3-235B} (Table~\ref{tab:qwen_pareto}) shows global $\rho{=}-0.12$ ($p{=}0.078$), with 13/20 languages exhibiting non-monotonic depth--transfer patterns.
Only Hindi reaches per-language significance ($\rho{=}-0.41$, $p{=}0.046$).
Final-layer transfer drop is 0.6pp, far below the pilot's 2--4pp.

\textbf{Llama~4 Scout} shows even weaker effects: global $\rho{=}-0.06$ ($p{=}0.31$), 0/20 per-language significant, 8/20 non-monotonic.
Final-layer drop is 0.3pp.

MoE expert routing breaks the depth$\to$specialization monotonic mapping that underlies Pareto in dense models.
Different experts are activated at different depths for different languages, creating non-monotonic transfer profiles.

\subsection{Dense Architecture: Marginal Pareto}

Gemma~4 31B Dense shows the clearest depth--transfer relationship: global $\rho{=}-0.24$ ($p{=}0.038$), 4/20 per-language significant (zh, fr, hi, sw), final-layer drop 0.8pp.
But even this is far weaker than the pilot: only 4/20 languages significant, and the 0.8pp final-layer drop is marginal.

\begin{table}[t]
\centering
\small
\caption{Dense vs MoE architecture comparison.}
\label{tab:arch_compare}
\begin{tabular}{lcccc}
\toprule
Metric & Qwen3 & Llama4 & Gemma4 \\
 & (MoE) & (MoE) & (Dense) \\
\midrule
EN best AUROC & .741 & .718 & .755 \\
Global $\rho$ & $-.12$ (n.s.) & $-.06$ (n.s.) & $-.24$ ($p$=.038) \\
Per-lang sig. & 1/20 & 0/20 & 4/20 \\
Non-monotonic & 13/20 & 8/20 & 5/20 \\
Final drop (pp) & 0.6 & 0.3 & 0.8 \\
ZS pass rate & 55\% & 45\% & 60\% \\
\bottomrule
\end{tabular}
\end{table}

Table~\ref{tab:arch_compare} summarizes: the Pareto tradeoff is architecture-conditional (present in dense, absent in MoE) and far weaker at full scale than the pilot predicted.
Three-architecture mean EN AUROC is 0.738.

\subsection{Transfer Conditions}

\begin{table}[t]
\centering
\small
\caption{Transfer conditions averaged across 3 architectures and 20 languages.}
\label{tab:transfer_5way}
\begin{tabular}{lcccc}
\toprule
Condition & AUROC & ECE & $\Delta$ AUROC & Win\% \\
\midrule
Zero-shot & .574 & .282 & --- & --- \\
Mean-shift & .574 & .288 & .000 & 0\% \\
CaCLR (best) & .566 & .275 & $-.008$ & 20\% \\
Lin.\ proj.\ (200) & .610 & .245 & $+.036$ & 70\% \\
Few-shot (25) & .596 & .258 & $+.022$ & 52\% \\
Few-shot (100) & .623 & .232 & $+.049$ & 75\% \\
\bottomrule
\end{tabular}
\end{table}

Table~\ref{tab:transfer_5way} shows the five-way comparison.
Mean-shift is provably inert (\S\ref{sec:meanshift}).
CaCLR degrades performance (\S\ref{sec:caclr}).
Linear projection and few-shot are the only beneficial conditions.

%-------------------------------------------------------------
\section{CKA Analysis: Simpson's Paradox Falsified}
\label{sec:cka}

The pilot found a Simpson's Paradox: raw CKA--transfer correlation was negative (depth confound), but partial correlation controlling for depth was positive and significant (Llama pilot: $\rho{=}+0.41$, $p{=}0.006$).
Full-scale results falsify this.

\begin{table}[t]
\centering
\small
\caption{Simpson's Paradox validation. Partial correlations controlling for depth are all non-significant.}
\label{tab:simpson}
\begin{tabular}{lccccl}
\toprule
Model & Raw $\rho$ & Raw $p$ & Partial $\rho$ & Partial $p$ & Verdict \\
\midrule
Qwen3 & $-.32$ & ${<}.001$ & $+.11$ & .28 & n.s. \\
Llama4 & $-.38$ & ${<}.001$ & $+.06$ & .45 & artifact \\
Gemma4 & $-.29$ & ${<}.001$ & $+.16$ & .09 & marginal \\
\bottomrule
\end{tabular}
\end{table}

Table~\ref{tab:simpson}: all three partial correlations are non-significant (max $\rho{=}+0.16$, min $p{=}0.09$).
The Llama pilot's $\rho{=}+0.41$ collapses to $+0.06$ with 20 languages, confirming it as a 5-language small-sample artifact.
CKA provides no explanatory power beyond depth.

\paragraph{Same-layer analysis.}
Within individual layers, CKA--transfer correlation is weak: median $\rho{=}0.14$ across all layers and models, with only 3/52 probed layers reaching significance.

\paragraph{Subspace angle analysis.}
Correctness direction angles between English and target languages vary dramatically by resource tier: HIGH 15--38\textdegree, MID 28--52\textdegree, LOW 42--65\textdegree\ (Table~\ref{tab:subspace}).
Within-language angle standard deviation across layers is 8--15\textdegree, indicating the correctness direction is neither language-agnostic nor layer-stable.

\begin{table}[t]
\centering
\small
\caption{Subspace angle between EN and target correctness directions at $\sim$60\% depth.}
\label{tab:subspace}
\begin{tabular}{llccc}
\toprule
Lang & Tier & Qwen3 & Llama4 & Gemma4 \\
\midrule
es & HIGH & 15\textdegree & 20\textdegree & 14\textdegree \\
zh & HIGH & 18\textdegree & 24\textdegree & 16\textdegree \\
ja & HIGH & 35\textdegree & 42\textdegree & 32\textdegree \\
ar & MID & 32\textdegree & 40\textdegree & 30\textdegree \\
ko & MID & 28\textdegree & 36\textdegree & 26\textdegree \\
hi & MID & 38\textdegree & 46\textdegree & 35\textdegree \\
tr & MID & 45\textdegree & 52\textdegree & 42\textdegree \\
sw & LOW & 52\textdegree & 58\textdegree & 48\textdegree \\
bn & LOW & 48\textdegree & 55\textdegree & 45\textdegree \\
am & LOW & 62\textdegree & 65\textdegree & 58\textdegree \\
\bottomrule
\end{tabular}
\end{table}

%-------------------------------------------------------------
\section{CaCLR and Distributional Matching}
\label{sec:caclr}

CaCLR optimizes $\mathcal{L} = \mathcal{L}_{\text{BCE}} + \alpha \cdot \mathcal{L}_{\text{SB-ECE}} + \beta \cdot \mathcal{L}_{\text{consistency}}$, where $\mathcal{L}_{\text{consistency}}$ uses MMD on probe logits from parallel inputs.
We test whether distributional matching can improve cross-lingual transfer.

\subsection{Implementation}
MMD uses a Gaussian RBF kernel with median-heuristic bandwidth.
Optimization: Adam, $\text{lr}{=}10^{-3}$ with cosine schedule to $10^{-5}$, 50 epochs, patience-10 early stopping on held-out English BCE loss.
Initialized from the standard LR solution.
Batch size 256.
Parallel set sizes range from 618 to 2,852 per language.

\subsection{$\beta$ Sweep: Consistency Loss is Inert or Harmful}

\begin{table}[t]
\centering
\small
\caption{CaCLR $\beta$ sweep (3 architectures $\times$ 20 languages).}
\label{tab:beta}
\begin{tabular}{lcccr}
\toprule
$\beta$ & AUROC & $\Delta$ vs ZS & \% Improved & \% Harmed \\
\midrule
0.001 & .572 & $-.002$ & 30\% & 42\% \\
0.01 & .569 & $-.005$ & 25\% & 52\% \\
0.1 & .565 & $-.009$ & 18\% & 62\% \\
1.0 & .558 & $-.016$ & 12\% & 75\% \\
\bottomrule
\end{tabular}
\end{table}

Table~\ref{tab:beta}: every $\beta$ value harms average performance.
Higher $\beta$ monotonically increases harm, with $\beta{=}1.0$ degrading 75\% of language--architecture combinations.
Multi-seed evaluation (5 seeds, std~$=$~0.003) confirms these are stable results.
A mixed $\beta$ criterion selecting on held-out zh+sw matches English-only selection ($\beta{=}0.001$), ruling out systematic bias against cross-lingual benefit.

\subsection{Distributional Matching: All Paths Closed}

We extend beyond 1-D MMD to test whether the failure is MMD-specific or fundamental.

\begin{table}[t]
\centering
\small
\caption{Distributional matching comprehensive ablation. All methods degrade performance. Linear projection (non-distributional) shown for contrast.}
\label{tab:distmatch}
\begin{tabular}{llc}
\toprule
Method & Dim. & $\Delta$ AUROC \\
\midrule
MMD (logit) & 1-D & $-.004$ \\
MMD (PCA-50) & 50-D & $-.006$ \\
MMD (PCA-100) & 100-D & $-.008$ \\
CORAL (logit) & 1-D & $-.005$ \\
CORAL (PCA-50) & 50-D & $-.007$ \\
\midrule
Linear projection & full-D & $+.036$ \\
\bottomrule
\end{tabular}
\end{table}

Table~\ref{tab:distmatch}: MMD fails across dimensionalities (1-D, 50-D, 100-D).
Higher dimensions are worse due to curse of dimensionality with ${\sim}2000$ training samples.
CORAL (second-order alignment) also fails.
This extends the negative result from ``MMD fails'' to ``distributional matching fails'': the cross-lingual correctness calibration problem violates the covariate shift assumption ($P(Y|X)$ differs across languages) that these methods require.

\subsection{Component Ablation}

\begin{table}[t]
\centering
\small
\caption{Component ablation: each added loss is net-negative.}
\label{tab:ablation}
\begin{tabular}{lcc}
\toprule
Configuration & $\Delta$ vs ZS & Note \\
\midrule
BCE only (retrain) & $-.001$ & $\approx$ ZS \\
BCE + SB-ECE & $-.005$ & ECE hurts discrim. \\
BCE + MMD (1-D) & $-.004$ & inert \\
BCE + MMD (PCA-50) & $-.006$ & fails \\
BCE + MMD (PCA-100) & $-.008$ & worse \\
BCE + CORAL & $-.005$ & fails \\
BCE + CORAL (PCA-50) & $-.007$ & fails \\
Full CaCLR & $-.008$ & compound degrad. \\
\bottomrule
\end{tabular}
\end{table}

Table~\ref{tab:ablation}: BCE-only retraining approximately recovers the zero-shot baseline (the LR solution is already optimal).
Each additional component degrades performance.
This is not a hyperparameter problem---it is a fundamental mismatch between the distributional matching objective and the conditional-distribution structure of cross-lingual correctness.

%-------------------------------------------------------------
\section{Difficulty Decomposition}
\label{sec:decomp}

We test whether transfer degradation decomposes into a difficulty-shared (language-agnostic) and capability-specific (language-dependent) component.
IRT difficulty probes ($R^2{=}0.28$) are trained, and the difficulty direction is projected out before retraining correctness probes.
A random-direction control projects out a randomly sampled direction of equal norm.

\begin{table}[t]
\centering
\small
\caption{Difficulty decomposition vs random-direction control (3-model mean). Paired $p$-values test decomposition vs random.}
\label{tab:decomp}
\begin{tabular}{llcccr}
\toprule
Lang & Tier & Orig.\ ZS & After proj. & $\Delta$ (pp) & Rand $\Delta$ \\
\midrule
zh & HIGH & .666 & .670 & +0.4 & +0.1 \\
es & HIGH & .674 & .676 & +0.2 & $-$0.1 \\
fr & HIGH & .657 & .660 & +0.3 & +0.2 \\
de & HIGH & .650 & .652 & +0.2 & $-$0.2 \\
hi & MID & .563 & .565 & +0.2 & +0.3 \\
ru & MID & .621 & .622 & +0.1 & $-$0.2 \\
ko & MID & .589 & .598 & +0.9 & +0.2 \\
tr & MID & .546 & .540 & $-$0.6 & $-$0.4 \\
sw & LOW & .515 & .510 & $-$0.5 & $-$0.3 \\
bn & LOW & .521 & .519 & $-$0.2 & $+$0.0 \\
\bottomrule
\end{tabular}
\end{table}

Mean decomposition improvement: +0.2pp (range: $-$0.8 to +1.2pp).
Random-direction control: $-$0.1pp (range: $-$1.0 to +0.8pp).
Paired test: $p{=}0.38$.
The decomposition is statistically indistinguishable from projecting out a random direction: observed deltas arise from dimensionality reduction effects, not difficulty-specific signal removal.
LOW-tier languages show negative deltas in both conditions---projecting out any direction damages their already-weak signal.
Difficulty and capability are entangled in the representation and cannot be operationally separated.

%-------------------------------------------------------------
\section{Control Task Analysis}
\label{sec:control}

We extend Hewitt \& Liang's control task framework to cross-lingual probing for the first time.
For each language--layer--model combination, we train probes on randomly permuted labels and compute selectivity $=$ AUROC$_{\text{real}}$ $-$ AUROC$_{\text{permuted}}$.
A combination is flagged as a potential artifact if selectivity $<$ 0.05.

\begin{table}[t]
\centering
\small
\caption{Cross-lingual artifact rates by resource tier (3 models, selectivity threshold 0.05).}
\label{tab:artifact}
\begin{tabular}{lcccc}
\toprule
Tier & $n$ langs & Artifact rate & Sel.\ mean & Sel.\ std \\
\midrule
HIGH & 7 & 15\% & .195 & .042 \\
MID & 7 & 31\% & .138 & .038 \\
LOW & 6 & 42\% & .082 & .035 \\
\midrule
ALL & 20 & 35\% & .142 & .058 \\
\bottomrule
\end{tabular}
\end{table}

\paragraph{35\% overall artifact rate.}
Table~\ref{tab:artifact}: more than a third of cross-lingual probe results fail the selectivity check.
LOW-resource languages suffer most (42\%), but MID-resource languages are also substantially affected (31\%).

\paragraph{Per-language selectivity.}
Table~\ref{tab:selectivity} shows per-language results for Qwen3.
Four LOW-resource languages (am, yo, my, sw) have selectivity $<$ 0.01 at the English-optimal layer---their probes learn capacity patterns, not correctness signal.
Three MID-resource languages (tr: 0.031, th: 0.028, vi: 0.042) also fall below or near the threshold.

\begin{table}[t]
\centering
\small
\caption{Per-language selectivity at EN-optimal layer (Qwen3-235B). ``$\dagger$'' indicates selectivity $<$ 0.05.}
\label{tab:selectivity}
\begin{tabular}{llcccc}
\toprule
Lang & Tier & Real & Perm. & Sel. & Flag \\
\midrule
en & --- & .741 & .488 & .253 & \\
zh & H & .668 & .495 & .173 & \\
es & H & .675 & .492 & .183 & \\
fr & H & .658 & .490 & .168 & \\
de & H & .649 & .486 & .163 & \\
ja & H & .588 & .501 & .087 & \\
pt & H & .662 & .489 & .173 & \\
it & H & .645 & .492 & .153 & \\
ar & M & .608 & .496 & .112 & \\
hi & M & .567 & .503 & .064 & \\
ru & M & .621 & .493 & .128 & \\
ko & M & .591 & .499 & .092 & \\
tr & M & .548 & .517 & .031 & $\dagger$ \\
vi & M & .562 & .520 & .042 & $\dagger$ \\
th & M & .528 & .500 & .028 & $\dagger$ \\
sw & L & .518 & .510 & .008 & $\dagger$ \\
bn & L & .524 & .505 & .019 & $\dagger$ \\
am & L & .486 & .482 & .004 & $\dagger$ \\
yo & L & .475 & .470 & .005 & $\dagger$ \\
te & L & .502 & .492 & .010 & $\dagger$ \\
my & L & .468 & .462 & .006 & $\dagger$ \\
\bottomrule
\end{tabular}
\end{table}

\paragraph{Selectivity predicts transfer.}
Selectivity correlates with zero-shot AUROC at $\rho{=}0.65$, suggesting that published positive cross-lingual probing results may partially reflect capacity artifacts rather than genuine correctness signal transfer.

\paragraph{Threshold sensitivity.}
At threshold 0.03, overall artifact rate rises to 42\% (LOW: 52\%); at 0.10, it drops to 24\% (LOW: 35\%).
The finding is robust across thresholds.

%-------------------------------------------------------------
\section{Linear Projection Analysis}
\label{sec:linproj}

Linear projection (Procrustes alignment) is the only consistently positive label-free method ($+$3.6pp vs zero-shot, Table~\ref{tab:transfer_5way}).
We analyze its sample efficiency and limitations.

\subsection{Sample Efficiency}

\begin{table}[t]
\centering
\small
\caption{Linear projection sample efficiency (3 architectures $\times$ 20 languages).}
\label{tab:sampleeff}
\begin{tabular}{lcccc}
\toprule
$n$ pairs & AUROC & \% of 200 & Std & Sig.\ \% \\
\midrule
10 & .552 & 66\% & .062 & 22\% \\
20 & .571 & 78\% & .055 & 38\% \\
50 & .592 & 89\% & .049 & 52\% \\
100 & .605 & 96\% & .044 & 65\% \\
200 & .610 & 100\% & .042 & 70\% \\
\bottomrule
\end{tabular}
\end{table}

Table~\ref{tab:sampleeff}: 100 pairs achieve 96\% of the 200-pair performance, but the pilot's claim that 50 pairs suffice is overly optimistic.
Saturation varies by tier: HIGH languages reach 95\% at $\sim$62 pairs (range 38--92), MID at $\sim$118 (72--158), LOW at $\sim$175 (140--198).
Low-resource languages require more pairs because their larger subspace angles (Table~\ref{tab:subspace}) and noisier translations demand more data to stabilize the Procrustes rotation.

\subsection{Effect Reversal}

\begin{table}[t]
\centering
\small
\caption{Procrustes effect by language (200 pairs). Effect reverses above ZS AUROC $\sim$0.62.}
\label{tab:reversal}
\begin{tabular}{llcccc}
\toprule
Lang & Tier & ZS & LinProj & $\Delta$ & Effect \\
\midrule
am & LOW & .483 & .508 & $+.025$ & helps \\
yo & LOW & .468 & .494 & $+.026$ & helps \\
sw & LOW & .511 & .528 & $+.017$ & helps \\
tr & MID & .544 & .560 & $+.016$ & helps \\
hi & MID & .560 & .568 & $+.008$ & marginal \\
vi & MID & .563 & .562 & $-.001$ & neutral \\
ko & MID & .589 & .582 & $-.007$ & hurts \\
ar & MID & .606 & .598 & $-.008$ & hurts \\
ru & MID & .622 & .610 & $-.012$ & hurts \\
es & HIGH & .674 & .654 & $-.020$ & hurts \\
ja & HIGH & .588 & .596 & $+.008$ & helps \\
\bottomrule
\end{tabular}
\end{table}

Table~\ref{tab:reversal}: below ZS AUROC $\sim$0.53, projection consistently helps; above $\sim$0.62, it consistently hurts.
The inconsistent zone (0.53--0.62) varies $\pm$0.08 across languages---Japanese (HIGH, 0.588) is helped while Korean (MID, 0.589) is hurt at nearly identical baseline AUROCs.
No single threshold is reliable.
Practical recommendation: apply projection when ZS AUROC $<$ 0.50 (likely beneficial), avoid when $>$ 0.65 (likely harmful), and validate on held-out data otherwise.

%-------------------------------------------------------------
\section{Supporting Results}

\subsection{Cross-Lingual Orthogonality}

Correctness--confidence orthogonality (cosine similarity) persists for HIGH-resource languages (mean cosine~$=$~0.035) but breaks down for LOW-resource (mean~$=$~0.108, std~$=$~0.078).
Architecture agreement is 65\% (13/20 languages).
The elevated LOW-resource cosine likely reflects probe noise rather than true non-orthogonality, consistent with the high artifact rates in \S\ref{sec:control}.

%-------------------------------------------------------------
\section{Discussion}
\label{sec:discussion}

\paragraph{Why do pilot results not generalize?}
The pilot used 7--8B dense models with 5 languages.
Dense architectures have monotonic depth$\to$specialization mappings; MoE routing breaks this.
Five languages over-represent HIGH-resource cases where effects are strongest.
Both the Pareto finding and the Simpson's Paradox were small-sample artifacts amplified by architectural specificity.

\paragraph{Why does distributional matching fail?}
Zhao et al. predict that marginal distribution matching fails when $P(Y|X)$ differs across domains.
Cross-lingual correctness probing violates covariate shift: $P(\text{correct}|\text{representation})$ differs systematically across languages because model capability varies.
MMD matches marginals but cannot align the conditional structure that carries correctness information.
This is confirmed by the monotonic degradation with $\beta$ (Table~\ref{tab:beta}) and the failure of both CORAL and high-dimensional MMD (Table~\ref{tab:distmatch}).

\paragraph{Field implications of control task findings.}
The 35\% artifact rate implies that cross-lingual probing studies---including those reporting only MID-resource languages---should routinely include selectivity checks.
The $\rho{=}0.65$ selectivity--AUROC correlation suggests that effect sizes in the literature may be inflated by capacity artifacts.

\paragraph{Limitations.}
(1)~MCQA-only evaluation; correctness probing may behave differently for open-ended generation.
(2)~Quantization artifacts at deep layers introduce 1--2pp bias for Qwen3 (INT4) and Llama4 (W4A16); Gemma4 (BF16) is unaffected (Appendix~\ref{app:quant}).
(3)~Linear projection requires parallel data; availability varies by language.

%-------------------------------------------------------------
\section{Conclusion}
\label{sec:conclusion}

We present the first systematic study of cross-lingual correctness probe transfer.
Our results close several paths:
distributional matching (MMD, CORAL, high-dimensional variants) consistently harms transfer;
mean-shift is provably inert;
the Pareto tradeoff is architecture-dependent and marginal at scale;
CKA provides no explanatory power beyond depth;
difficulty--capability decomposition is not operational.

Two positive contributions remain.
First, control task analysis reveals that 35\% of cross-lingual probe results are capacity artifacts---a methodological warning for the field.
Second, linear projection with ${\geq}100$ parallel pairs is the only consistently beneficial label-free method, though its effect reverses for high-performing languages.
These findings prevent the community from pursuing demonstrably closed paths while providing conditional practical guidance.

%-------------------------------------------------------------
\appendix

\section{Probe Complexity Comparison}
\label{app:probe}

\begin{table}[h]
\centering
\small
\caption{Logistic regression vs 2-layer MLP (3 models $\times$ 20 languages).}
\label{tab:probecmp}
\begin{tabular}{lcccc}
\toprule
Metric & LR & MLP & $\Delta$ & $p$ \\
\midrule
EN AUROC & .738 & .752 & +.014 & .006 \\
ZS degradation & 15.8\% & 19.4\% & +3.6pp & .018 \\
ZS AUROC & .621 & .608 & $-.013$ & .032 \\
EN ECE & .238 & .252 & +.014 & .058 \\
LR better for xfer & 72\% & 28\% & --- & --- \\
\bottomrule
\end{tabular}
\end{table}

MLP gains +1.4pp EN AUROC but transfers worse ($-$3.6pp more degradation).
LR is preferred for cross-lingual transfer in 72\% of languages.

\section{Cross-Benchmark Generalization}
\label{app:benchmarks}

\begin{table}[h]
\centering
\small
\caption{Cross-benchmark generalization.}
\label{tab:benchmark}
\begin{tabular}{lccccc}
\toprule
Benchmark & ZS pass & ZS AUROC & CaCLR $\Delta$ & LinProj $\Delta$ & Corr. \\
\midrule
GlobalMMLU & 57\% & .574 & $-.008$ & $+.036$ & --- \\
Belebele & 45\% & .545 & $-.015$ & $+.028$ & .71 \\
MMLU-ProX & 55\% & .565 & $-.009$ & $+.033$ & .79 \\
\bottomrule
\end{tabular}
\end{table}

Cross-benchmark correlation is moderate (.71--.79).
Belebele (reading comprehension) shows lower transfer, suggesting correctness probing is partially MCQA-format-specific.
CaCLR negative results replicate across benchmarks.

\section{Quantization Validation}
\label{app:quant}

\begin{table}[h]
\centering
\small
\caption{Quantization impact on probe AUROC (5 representative layers).}
\label{tab:quant}
\begin{tabular}{llcccl}
\toprule
Model & Quant & Depth\% & FP16 & Quant & $\Delta$ \\
\midrule
Llama4 & W4A16 & 20\% & .650 & .647 & $-.003$ \\
Llama4 & W4A16 & 60\% & .698 & .688 & $-.010$ \\
Llama4 & W4A16 & 95\% & .715 & .699 & $-.016$ \\
\midrule
Qwen3 & INT4 & 20\% & .628 & .623 & $-.005$ \\
Qwen3 & INT4 & 60\% & .705 & .692 & $-.013$ \\
Qwen3 & INT4 & 95\% & .741 & .720 & $-.021$ \\
\bottomrule
\end{tabular}
\end{table}

Quantization introduces 0.3--2.1pp AUROC degradation, increasing with depth.
Rank preservation: Llama4 92\%, Qwen3 89\%.
Deep-layer conclusions are potentially affected, but the negative results (distributional matching failure, decomposition null) do not depend on precise values.
Gemma4 (BF16) provides quantization-free validation.

\section{Extended Language Coverage}
\label{app:extended}

Extending from 20 to 52 languages: ZS pass rate drops from 57\% to 50\% (more LOW-resource languages added).
LOW-tier artifact rate rises to 45\%.
Core conclusions are unchanged; expanded coverage strengthens the control task finding.

% Bibliography omitted for draft
%\bibliography{refs}
%\bibliographystyle{plainnat}

\end{document}
